% !TEX root = ../../../main.tex
% 来源：中学数学实验教材·第二册（几何）
% 章节：圆 / 圆与直线的位置关系
% 原始文件：4.tex，行 1876-1896
% 类型：tikz
% ---
\begin{tikzpicture}
\tkzDefPoints{0/0/I, 0/-1.5/Q}
\tkzDefPoint(10:1.5){R}
\tkzDefPoint(165:1.5){P}
\tkzDefPoint(75:1.5){S}
\tkzDefTangent[at=R](I) \tkzGetPoint{a}
\tkzDefTangent[at=S](I) \tkzGetPoint{b}
\tkzDefTangent[at=P](I) \tkzGetPoint{c}
\tkzDefTangent[at=Q](I) \tkzGetPoint{d}
\tkzInterLL(R,a)(S,b)\tkzGetPoint{C}
\tkzInterLL(S,b)(P,c)\tkzGetPoint{D}
\tkzInterLL(P,c)(Q,d)\tkzGetPoint{A}
\tkzInterLL(R,a)(Q,d)\tkzGetPoint{B}
\tkzDrawPolygon[thick](A,B,C,D)
\tkzDrawCircle[thick](I,Q)
\tkzLabelPoints[below](A,Q,B)
\tkzLabelPoints[left](D,P)
\tkzLabelPoints[right](C,R,I)
\tkzLabelPoints[above](S)
\tkzDrawPoints(I)
\end{tikzpicture}
